% Lösungen Einheit 3 von 4 – Gerade und Ebene (zusammenbau v0.1)
\einheitenkopf{Einheit 3 von 4 · Gerade und Ebene}

\erg{14}{a) liegt in der Ebene – $t$ fällt heraus, und $4 = 4$ stimmt. \quad b) $(1 + t) + 2(2 - t) + t = 5$, also $5 = 5$ für jedes $t$ – $g$ liegt in $E$. \quad c) $\vec{u} \cdot \vec{n} = 2 - 2 + 1 = 1 \neq 0$ – genau ein Schnittpunkt. \quad d) nein – die z-Koordinaten $8$ und $6$ hängen nicht von $t$ ab; der Richtungsvektor hat für jedes $t$ die z-Komponente $-2$. \quad e) Der Normalenvektor $(1 | 2 | 0)$ von $L$ ist der Richtungsvektor der Achse, also steht $L$ senkrecht auf $a$; $S$ liegt in $L$, denn $4 + 2 \cdot 3 = 10$; bei der Drehung bleibt $S$ in dieser Ebene, der Mittelpunkt liegt auf der Achse: $t + 4t = 10$, $t = 2$, also $H(2 | 4 | 2)$. \quad f) Die Punkte $(a | a | 0)$ bilden die Gerade durch den Ursprung mit dem Richtungsvektor $(1 | 1 | 0)$. Jede Ebene $x - y + z = c$ mit $c \neq 0$ ist echt parallel zu ihr: $1 - 1 + 0 = 0$, und der Ursprung erfüllt die Gleichung nicht; $c$ ist beliebig. \quad g) $\vec{u} \cdot \vec{n} = 2a - 2$. Für $a = 1$ ist es null, und der Stützpunkt $(1 | 0 | 1)$ erfüllt $2 - 1 = 1$: $g_1$ liegt in $E$. Für $a \neq 1$: $2(1 + ra) - (a + 2r) = 1$ gibt $r = \frac{1}{2}$, Schnittpunkt $S_a(1 + \frac{a}{2} | 1 | a + 1)$.}
\erg{15}{a) Nils hat nur den Stützpunkt geprüft; eingesetzt: $2(1 + t) - 2 + t = 3t = 0$ nur für $t = 0$ – $g$ schneidet $E$ im Stützpunkt. \quad b) Die Gleichung gilt dann für jeden Parameterwert, also für jeden Punkt der Geraden.}
